% Seção «Inteligência Artificial e Decisões Automatizadas», versão LaTeX (abnTeX2).
% Gerada do corpo de docs/secao-ia-decisoes-automatizadas.md (sha256 0ca45d4c8e28916e) por
% scripts/gera_secao_ia_latex.py, com pandoc 3.1.3.
% NÃO EDITE AQUI: edite o .md e gere de novo. O .md é a fonte, e a conferência
% de cada número está nele. Para o Overleaf: colar este arquivo inteiro no
% capítulo escolhido (ele começa em um \section); as
% 20 referências novas estão em secao-ia-decisoes-automatizadas.bib, e as
% outras 4 citadas já estão no .bib do Overleaf. Requer a classe abntex2
% (\fonte e \nota) e o abntex2cite (\cite); nenhum pacote a mais.

% A classe abntex2 traz \fonte e \nota, mas não o ambiente quadro: quem o define
% é o preâmbulo do modelo do abnTeX2. Se o do Overleaf já o tiver, as linhas abaixo
% não fazem nada.
\makeatletter
\@ifundefined{quadro}{%
  \providecommand{\quadroname}{Quadro}%
  \newfloat{quadro}{loq}{\quadroname}%
}{}
\makeatother
% As tabelas e o quadro desta seção não flutuam: ficam no lugar do texto, sem cortar
% parágrafo. A legenda fixa é a da memoir, a base do abntex2, e numera e entra na lista
% como a legenda de sempre.
\providefixedcaption{\tabelafixa}{table}
\providefixedcaption{\quadrofixo}{quadro}
\section{Inteligência Artificial e Decisões Automatizadas}

Na ferramenta desenvolvida, algoritmos de inteligência artificial apoiam as decisões metodológicas do refino quimiométrico, e esse apoio se concentra nas escolhas que antecedem o processamento de uma corrida analítica: a forma do erro sistemático presente nos dados, o número de lotes analíticos a declarar e o método de correção de sinal a aplicar. Declarar lotes significa dividir a corrida em trechos, cada um corrigido com seus próprios QCs.

Escolher, a partir de características mensuráveis de uma instância, o algoritmo com o melhor desempenho esperado sobre ela é o chamado problema de seleção de algoritmos \cite{rice1976algorithm}. Aprender essa escolha com base em experiências anteriores é o objeto do meta-aprendizado, uma das bases do aprendizado de máquina automatizado (AutoML) \cite{brazdil2022metalearning}. Em metabolômica, ferramentas como o NormalizeMets comparam métodos de normalização sobre os dados do próprio usuário \cite{de2018normalizemets}, e o aprendizado de máquina tem sido aplicado à modelagem de doenças e à classificação de amostras \cite{galal2022applications}.

Na ferramenta, o apoio à decisão combina mecanismos de naturezas diferentes. O diagnóstico do erro sistemático é um classificador baseado em regras, construído e calibrado neste trabalho. A escolha do corretor, nome dado aqui ao método de correção de sinal, foi estudada primeiro como um modelo de aprendizado de máquina e acabou resolvida por medição, como se descreve adiante. O aprendizado de máquina aparece ainda dentro dos próprios métodos de correção e de análise.

O princípio que orientou o desenho desses pontos é que a decisão automática funcione como uma segunda opinião. As propostas chegam ao usuário com a justificativa e podem ser aceitas ou alteradas, e cada execução fica registrada com a escolha efetivamente usada. Esse cuidado se justifica porque auxílios à decisão imperfeitos levam a dois tipos de erro: no de omissão, o usuário deixa de notar um problema que o auxílio não apontou, e no de comissão, segue uma recomendação errada sem verificá-la. Os dois caracterizam o viés de automação, que não é evitado por treinamento nem por instruções \cite{parasuraman2010complacency}. O Quadro~\ref{quadro:ia-decisoes} resume os pontos em que o aplicativo decide ou propõe algo no lugar do usuário.

\par\medskip
\noindent\begin{minipage}{\linewidth}
  \quadrofixo{Pontos de decisão automática do aplicativo}
  \label{quadro:ia-decisoes}
  \centering
  \footnotesize
  \setlength{\tabcolsep}{3pt}
  \begin{tabular}{|>{\raggedright\arraybackslash}p{1.9cm}|>{\raggedright\arraybackslash}p{3.0cm}|>{\raggedright\arraybackslash}p{3.1cm}|>{\raggedright\arraybackslash}p{3.0cm}|>{\raggedright\arraybackslash}p{2.7cm}|}
    \hline
    \textbf{Ponto de decisão} & \textbf{O que a máquina propõe ou decide} & \textbf{Mecanismo} & \textbf{O que o usuário vê} & \textbf{Controle do usuário} \\
    \hline
    Diagnóstico & a classe do erro sistemático e as fronteiras de lote; o número de lotes é pré-preenchido & sequência ordenada de regras sobre o RSD dos QCs e sobre os componentes principais, com detecção de pontos de mudança na trajetória dos QCs & classe, confiança, pontos de corte e número de QCs em cada lote proposto & altera o número de lotes; mantido o proposto, a correção corta nas fronteiras propostas \\
    \hline
    Palpite (a sugestão de corretor) & o corretor (não corrigir, LOESS ou SERRF) e a confiança na escolha, ou o empate & a mesma validação da linha seguinte, em 200 \emph{features} sorteadas da matriz & corretor, confiança e justificativa & caixa ``Automático'', marcada por padrão, ou escolha manual \\
    \hline
    Medir qual corretor vence & a opção de menor RSD dos QCs fora da amostra, inclusive não corrigir, ou a abstenção & validação que retira um QC por vez, nas três opções, com os lotes declarados & RSD de cada opção, margem, abstenção e se o palpite coincidiu & nenhum: o resultado é exibido e não altera os controles \\
    \hline
    Teste univariado & teste \emph{t} de Welch ou teste de Mann-Whitney & teste de normalidade de Shapiro-Wilk sobre os desvios dentro dos grupos & o teste escolhido e o \emph{p}-valor do teste de normalidade, no gráfico de vulcão & a janela não permite escolher o teste \\
    \hline
  \end{tabular}
  \fonte{elaborado pelos autores.}
\end{minipage}
\par\medskip

\subsection{Diagnóstico automático do erro sistemático}

O estudo usado neste trabalho tem quatro corridas, uma por modo analítico. Cada modo combina uma coluna cromatográfica (HILIC, de interação hidrofílica, ou C18, de fase reversa) com uma polaridade de ionização, positiva ou negativa. A análise exploratória mostrou que o erro sistemático não tem a mesma forma nas quatro. Projetados no primeiro componente principal, os QCs do HILIC negativo mudam de nível entre as injeções 45 e 56 e permanecem deslocados, os do C18 negativo se deslocam de forma gradual e os dos dois modos positivos oscilam sem padrão monotônico. Os métodos de correção partem de premissas diferentes sobre essa forma. A regressão local (LOESS) ao longo da ordem de injeção supõe uma deriva suave, e um degrau só é corrigido se as duas partes da corrida forem declaradas como lotes distintos. Por isso, a ferramenta classifica o tipo de defeito antes de propor a correção.

A diretriz de consenso de Broadhurst e colaboradores recomenda inspecionar o RSD de cada metabólito nos QCs, a posição dos QCs no gráfico de escores dos dois primeiros componentes principais e a resposta dos metabólitos ao longo da ordem de injeção \cite{broadhurst2018guidelines}. As recomendações do QComics acrescentam cartas de controle com as regras de Westgard, que detectam medições fora de controle \cite{doi:10.1021/acs.analchem.3c03660}. Essas regras indicam se a corrida saiu de controle, mas não qual é a forma do erro, e, com sete QCs por corrida, a que exige dez pontos consecutivos do mesmo lado da média nem chega a ser aplicável. A distinção entre formas de erro tem precedente em proteômica, em que se separam efeitos de lote contínuos, corrigidos por regressão ao longo da ordem das corridas, de efeitos discretos, corrigidos por centragem \cite{cuklina2021diagnostics}. Neste trabalho, essa separação é aplicada à matriz inteira por uma sequência de regras com limiares explícitos.

O diagnóstico transforma as inspeções da diretriz em quatro grandezas, calculadas sobre a matriz ainda não corrigida. A primeira é a mediana, entre as \emph{features}, do RSD dos QCs, calculado sobre as intensidades sem transformação e com a mesma remoção de valores discrepantes que o pipeline do laboratório aplica. As outras três vêm da análise de componentes principais da matriz transformada por \(\log_{10}(1 + x)\) e autoescalonada: a correlação de Spearman, em valor absoluto, entre o primeiro componente e a ordem de injeção; a razão entre a dispersão dos QCs e a das amostras nos dois primeiros componentes; e a trajetória dos QCs no primeiro componente. A agregação pela mediana e a correlação com a ordem de injeção são resumos construídos neste trabalho, já que a diretriz de Broadhurst e colaboradores descreve as inspeções metabólito a metabólito ou por gráfico. A correlação com a ordem de injeção e a razão de dispersão são ainda recalculadas sob pequenas perturbações dos escores, com semente fixa, e uma métrica indefinida ou instável não é usada para classificar.

As regras são avaliadas em ordem fixa, e a primeira que se aplica decide a classe:

\begin{enumerate}
\item
  \emph{não avaliável}, quando as métricas de componentes principais são indefinidas ou instáveis, ou quando a trajetória dos QCs quase não se move, o que corresponde a uma abstenção do diagnóstico;
\item
  \emph{limpo}, quando a mediana do RSD dos QCs fica abaixo do limiar do modo e a correlação do primeiro componente com a ordem de injeção fica abaixo de 0,30;
\item
  \emph{mudança em degrau}, quando há ponto de mudança na trajetória dos QCs e o novo nível se sustenta, isto é, quando a diferença entre as médias dos dois segmentos é mais que o dobro da maior amplitude interna a eles;
\item
  \emph{deriva gradual}, quando a trajetória é predominantemente monotônica, com razão de 0,487 ou mais entre o deslocamento líquido e o deslocamento total;
\item
  \emph{errático}, nos demais casos.
\end{enumerate}

Cada classificação traz uma confiança, calculada pela distância relativa entre o valor medido e o limiar que decidiu. Uma decisão tomada perto do limiar recebe confiança baixa, e esse valor não é uma probabilidade. O limiar de RSD da regra de limpeza é calibrado por modo, sobre matrizes sintéticas, com a restrição de manter nas corridas reais as classes vistas na análise exploratória: ele fica em cerca de 105\% no HILIC positivo, 23\% no C18 positivo, 49\% no HILIC negativo e 16\% no C18 negativo. Como acompanha o nível de RSD de cada modo, a classe limpo indica que a corrida não se afasta do nível do seu modo nem acompanha a ordem de injeção, e não que sua variabilidade seja baixa. Para um modo que a ferramenta não conhece, vale o corte de 20\%, o mais estrito dos dois, 20\% e 30\%, que a diretriz cita \cite{broadhurst2018guidelines}. Os limiares foram escolhidos a partir de medições sobre as quatro corridas e são próprios desses dados, e não constantes gerais do método. Com sete QCs, a regra da deriva gradual tem ainda um limite estatístico: em simulação, 36\% das trajetórias de um passeio aleatório puro, lidas nas posições dos sete QCs, já passam de 0,487, e subir o limiar não separaria a deriva do passeio aleatório.

Para propor fronteiras de lote, o diagnóstico procura pontos de mudança na trajetória dos QCs por segmentação binária, com um modelo constante em cada segmento \cite{truong2020selective}. Um corte é aceito quando reduz a soma dos quadrados dos resíduos em mais que uma penalidade de vinte vezes a variância do ruído. O ruído é estimado pelas diferenças entre QCs consecutivos, que um degrau quase não altera. O fator 20 fica perto do centro, em escala logarítmica, da faixa de 2,5 a 200 em que o detector acertou os casos de teste construídos e manteve um único corte no HILIC negativo. Como a trajetória tem um ponto por QC, cada corte é uma hipótese sobre poucos pontos, e a janela o apresenta como sugestão, com o número de QCs de cada lote proposto. Nas corridas reais, o detector propõe dois lotes no HILIC negativo, com o corte antes da injeção 56, dois no C18 negativo, um no C18 positivo e três no HILIC positivo. A busca de fronteiras corre independentemente da classe do defeito. Quando o usuário mantém o número proposto, a correção divide a corrida nessas fronteiras; se ele o altera, a corrida é dividida em blocos contíguos de tamanho quase igual. A diferença pode ser grande: no HILIC negativo, dois blocos iguais separam as injeções 34 e 35, e o degrau ficaria dentro do segundo lote.

Também foi testado se declarar dois lotes no HILIC negativo, com a fronteira logo depois da injeção 45, o último QC antes do degrau, reduz a correlação do primeiro componente com a ordem de injeção depois do SERRF. A correlação caiu de 0,61, com um lote, para 0,18, com dois. Em vários outros pontos de corte, porém, a redução foi igual ou maior: dos outros 43 cortes que deixam ao menos dois QCs em cada lote, 17 deram correlação menor, a maioria com o segundo lote começando perto do meio da corrida. A redução, portanto, não mostra que o degrau foi corrigido. A menor correlação de todas, 0,08, veio dos cortes antes das injeções 55 e 56, o segundo deles o proposto pelo detector. Nesses cortes, o segundo componente passa a ter correlação de 0,84 a 0,92 com a ordem de injeção, e a ordem continua pesando nos dados.

Nas corridas reais, o diagnóstico classifica o HILIC negativo como mudança em degrau, o C18 negativo como limpo e os dois modos positivos como erráticos. O C18 negativo sai como limpo, apesar da deriva vista na análise exploratória, porque o RSD mediano dos seus QCs fica abaixo do limiar do modo e o primeiro componente quase não acompanha a ordem de injeção. Essas classificações não medem o desempenho do método, porque os limiares foram escolhidos para reproduzi-las e a forma verdadeira do erro nessas corridas não é conhecida. O desempenho foi medido em matrizes produzidas pelo gerador de dados sintéticos desenvolvido neste trabalho, que reamostra os resíduos das corridas reais e injeta um defeito de forma e intensidade conhecidas, o que permite avaliar o método contra uma resposta conhecida de antemão \cite{morris2019using}. A intensidade é dada pela razão entre a amplitude da deriva e a variação de fundo de cada \emph{feature}.

Em 200 matrizes com sete QCs, 50 de cada classe, e razão de 0,8, a classe correta foi recuperada em 139 (69,5\%), contra 25\% esperados ao acaso. O acerto é alto nas matrizes limpas (46 de 50) e mais baixo na mudança em degrau (26 de 50). Ele depende da intensidade do defeito: com razões de 0,2, 0,4, 1,5 e 3,0, foi de 35,0\%, 46,0\%, 88,0\% e 88,5\%. Com deriva fraca, a maior parte das matrizes com defeito é classificada como limpa ou errática; com deriva forte, o degrau e a deriva gradual são reconhecidos em todas as matrizes, e os erros se concentram nas erráticas. O acerto também varia com o modo de origem, de 50,0\% nas matrizes derivadas do HILIC positivo a 91,7\% nas do C18 negativo, e cai para 58,0\% quando se usa o limiar de reserva de 20\%, o de um modo desconhecido. A confiança declarada acompanha o acerto: nas 800 matrizes distintas das cinco razões (as 50 limpas são comuns a todas), a classe está certa em 29,5\% das vezes com confiança abaixo de 0,1 e em 69,6\% com confiança a partir de 0,3.

A proposta de lotes foi avaliada nas mesmas matrizes. Numa matriz em degrau, o corte certo é o que inicia um lote no primeiro QC posterior ao degrau. Nas 50 matrizes em degrau de cada razão, da 0,2 à 3,0, o detector propôs só o corte certo em 4, 7, 31, 46 e 47 delas. Como independe da classe, a proposta também divide matrizes sem degrau: recebem algum corte 9 das 50 limpas e, na razão 0,8, 35 das 50 de deriva gradual. Daí a janela mostrar o número de QCs de cada lote e deixar o número de lotes sob controle do usuário.

O diagnóstico não escolhe o corretor. Nas 88 matrizes sintéticas usadas na escolha do corretor, descrita a seguir, todas com deriva gradual injetada, o SERRF venceu pela medida dos QCs em todas as classes que o diagnóstico atribuiu: em 15 das 17 classificadas como limpas, 22 das 24 erráticas, 35 das 46 de deriva gradual e na única em degrau. O tipo de defeito, sozinho, não indica o corretor, e a ferramenta não usa uma regra que associe um ao outro.

\subsection{Escolha automática do corretor}

A segunda decisão automática é a escolha entre as três opções de correção de sinal. A primeira é não corrigir, que é uma opção legítima, porque a correção também pode piorar o dado. A segunda é o LOESS, ajustado aos QCs ao longo da ordem de injeção, na linha da correção proposta por Dunn e colaboradores \cite{dunn2011procedures}. A terceira é o SERRF \cite{fan2019systematic}, que modela o erro de cada \emph{feature} com uma floresta aleatória treinada nos QCs. Das etapas do pré-processamento, só o corretor é escolhido automaticamente; as demais seguem os padrões declarados da ferramenta ou a escolha do usuário.

No fluxo anterior do laboratório, a correção roda no pipeline em R, que aplica o SERRF e o LOESS, e a escolha entre eles é feita olhando o resultado. Segundo a equipe, espera-se que, depois de uma boa correção, os QCs apareçam agrupados e no centro do gráfico de componentes principais, e escolhe-se o corretor com que os QCs ficaram mais agrupados e com menor RSD. O pipeline grava, junto das matrizes corrigidas, um gráfico de barras do RSD dos QCs e o gráfico de componentes principais. O RSD das barras é medido em dois QCs retirados do ajuste, em cinco sorteios sem semente, e por isso muda de uma execução para outra. O agrupamento no gráfico, ao contrário, avalia a correção nos próprios QCs que a ajustaram, e o LOESS, que ajusta uma curva aos QCs, reduz a variação justamente neles. O artigo do SERRF já advertia que o desempenho devia ser avaliado em amostras que não participaram do treinamento \cite{fan2019systematic}.

A ferramenta adota uma medida fora da amostra que se repete. Cada QC é retirado por vez, o corretor é reajustado com os demais e o RSD é calculado sobre os valores corrigidos dos QCs retirados. O corretor de menor mediana desse RSD, entre as \emph{features}, é o vencedor, e quando a segunda melhor opção fica a menos de 1,2\% da melhor, em termos relativos, a medição se abstém. Esse limiar vem do ruído da própria medida. Repetida com 21 sementes diferentes nos quatro arquivos reais, a margem entre as duas melhores opções teve desvio-padrão de 0,20\% a 0,51\%, conforme o arquivo, e em nenhum deles a troca de semente mudou o vencedor. O dobro do maior desses desvios dá 1,03\%, com intervalo de 95\% por reamostragem de 0,82\% a 1,24\%, e o limiar de 1,2\% fica dentro dele: uma margem abaixo disso não se distingue do que a troca de semente produz sozinha. No C18 negativo, onde a margem é a menor dos quatro, uma das 21 sementes a deixou abaixo do limiar, e a medição teria se abstido.

A primeira versão da escolha automática tentava prever esse vencedor sem corrigir a matriz, como no problema de seleção de algoritmos. Era uma política aprendida, treinada em matrizes sintéticas, porque o estudo tem só quatro corridas reais. O conjunto usado para avaliá-la reúne 88 matrizes, 22 derivadas de cada modo, com 2.000 \emph{features}, sete QCs nas posições das corridas reais e razão entre deriva e ruído sorteada, em escala logarítmica, entre 0,1 e 3,0. Cada matriz recebe como rótulo o vencedor da medida fora da amostra, o que exige reajustar os três corretores uma vez para cada QC retirado.

As entradas da política são 17 características que o diagnóstico calcula da matriz bruta, como a mediana e o percentil 75 do RSD dos QCs, a fração de zeros, a correlação do primeiro componente com a ordem de injeção, a razão de dispersão dos QCs, a confiança e o tipo de defeito. Ficam de fora, de propósito, as grandezas que dependem de executar os corretores. Se a política pudesse lê-las, a resposta já seria conhecida antes da recomendação, e a avaliação seria contaminada por informação indisponível no momento da decisão, o que se conhece como vazamento \cite{kaufman2012leakage}. O modelo é formado por duas árvores de decisão do tipo CART \cite{breiman1984classification}: uma classifica o corretor, e a outra prediz a margem entre o melhor e o segundo, de onde sai a confiança. A profundidade foi limitada a três níveis para que o caminho da raiz até a folha, com a característica, o valor medido e o limiar de cada nó, sirva de justificativa da decisão, como recomenda o uso de modelos interpretáveis quando a decisão tem consequências \cite{rudin2019stop}.

A política foi comparada com a regra mais simples possível, escolher sempre o SERRF, o corretor que vence em 73 das 88 matrizes. Como as duas julgam as mesmas matrizes, a diferença foi avaliada pelo teste exato de McNemar sobre os casos em que discordam \cite{mcnemar1947note}, indicado para comparar dois classificadores no mesmo conjunto de teste \cite{dietterich1998approximate}. O desempenho foi estimado de três formas. Avaliada nas próprias matrizes do treinamento, a política acerta 78. Retirando uma matriz por vez, acerta 65, menos que a regra fixa, com 10 discordâncias contra e 2 a favor (\emph{p} = 0,04). Retirando todas as matrizes de um modo e julgando esse modo com o modelo treinado nos outros três, acerta 72, contra 73, com uma única discordância. Esta última forma é a que importa, porque matrizes do mesmo modo compartilham a estrutura de correlação, o padrão de zeros e o perfil dos QCs, e separar treinamento e teste por matriz deixa essa informação dos dois lados \cite{roberts2017cross}. A política e a regra fixa quase não discordam nessa forma de avaliação, e nenhuma versão, nesta nem nas treinadas antes, mostrou vantagem estatisticamente detectável sobre ela. A política aprendida continua disponível pela linha de comando, mas deixou de ser a sugestão do aplicativo.

O palpite, nome que o aplicativo dá à sugestão de corretor, passou então a medir. O aplicativo aplica as três opções a 200 \emph{features} sorteadas da matriz do usuário, com semente fixa, e propõe a de menor RSD dos QCs fora da amostra. Ele continua sendo um palpite, porque mede uma amostra e não todas as \emph{features}, mas agora apoiado numa medição. A escolha de 200 \emph{features} considerou o custo, que cresce linearmente com o número de \emph{features}, e uma simulação sobre os valores por \emph{feature} já medidos nas 88 matrizes sintéticas, que têm todas deriva gradual. Refeita com o código atual, com 200 sorteios por matriz, a concordância com a medição de todas as \emph{features} foi de 86,1\% com 50 \emph{features}, 89,6\% com 100, 93,1\% com 200 e 95,6\% com 500, enquanto escolher sempre o SERRF acertaria 73 das 88 matrizes (83,0\%). Com 200, quando a amostra erra, o RSD da opção escolhida fica em média cerca de 4\% acima do da melhor, em diferença relativa; somados todos os sorteios, a diferença média é de 0,28\%.

Quando a diferença relativa entre as duas melhores opções na amostra fica abaixo de cerca de 17\%, o dobro do maior desvio-padrão dessa diferença entre sorteios nos quatro arquivos reais, o palpite declara empate, mas mantém a proposta. Esse limiar é bem mais largo que o de 1,2\% da medição de todas as \emph{features}, porque, com 200 \emph{features}, a diferença muda muito mais de um sorteio para outro do que com a troca de semente. Substituir a proposta pelo SERRF nos empates foi avaliado, num outro conjunto de sorteios, e aumentou a diferença média para a melhor opção de 0,26\% para 0,78\%. Nos quatro arquivos reais, o palpite propôs o mesmo corretor que a medição de todas as \emph{features}, e em dois deles, o HILIC negativo e o C18 negativo, declarou empate. Num computador de 12 núcleos lógicos, ele levou menos de um minuto por arquivo.

A caixa ``Automático'' do corretor vem marcada e entrega a escolha ao palpite no momento da execução. Quem executa sem alterar nada recebe o corretor do palpite, e nesse uso a decisão automática deixa de ser segunda opinião e passa a ser a escolha. Por isso, ao fim da execução, a janela informa qual corretor foi aplicado e quem o escolheu. Para quem quer o resultado completo, o botão ``Medir qual corretor vence'' aplica a mesma medida a todas as \emph{features}, com os lotes declarados. No mesmo computador, com um lote, essa medida levou de 8,6 minutos (HILIC negativo) a 27 minutos (HILIC positivo). Como cada QC retirado precisa de pelo menos dois outros no mesmo lote, ela recusa uma divisão com menos de três QCs em algum lote, o que, com sete QCs, admite no máximo dois lotes. A proposta do detector não pôde ser medida no HILIC negativo (5 e 2 QCs) nem no HILIC positivo (2, 2 e 3). Nas divisões em dois lotes que puderam ser medidas, no C18 negativo e em blocos iguais no HILIC negativo, nenhuma superou o lote único.

Cada execução fica registrada no banco de dados local do aplicativo, com a configuração efetivamente usada e a justificativa da escolha. Quando o usuário altera um valor proposto, o registro guarda também a proposta da máquina e quais parâmetros foram alterados, e a substituição de uma proposta pode ser contada depois. Manter a proposta sem mexer no controle, porém, não se distingue de avaliá-la e mantê-la, e nenhuma sessão de uso por especialista foi registrada até aqui.

Nos quatro arquivos reais, com um lote, a medição elegeu o SERRF no HILIC positivo, no C18 positivo e no C18 negativo, e o LOESS no HILIC negativo. As medianas estão na Tabela~\ref{tab:ia-reais}, ao lado das do pipeline do laboratório.

\subsection{O SERRF e o LOESS da ferramenta e os do laboratório}

O SERRF da ferramenta é uma reimplementação do original em R. A floresta aleatória \cite{breiman2001random} de cada \emph{feature} é treinada apenas nos QCs e prevê a intensidade a partir das dez \emph{features} mais correlacionadas com ela \cite{fan2019systematic}. As duas implementações diferem em reprodutibilidade. O SERRF do pipeline em R não fixa semente e substitui os zeros por valores sorteados, de modo que duas execuções sobre o mesmo arquivo não dão o mesmo resultado. O LOESS do pipeline também tem partes sorteadas sem semente, e, quando mais de um quinto da curva ajustada de uma \emph{feature} sai negativa, devolve a \emph{feature} sem correção e sem avisar, o que, numa execução, aconteceu em 1.561 das 6.035 \emph{features} do HILIC positivo (25,9\%). Na ferramenta, toda operação aleatória recebe uma semente derivada da semente geral e da \emph{feature} (no SERRF, também do lote e do passo), a mesma entrada dá sempre a mesma saída, e as \emph{features} que o LOESS deixa sem correção ficam marcadas.

A diferença tem consequência prática no fluxo anterior, em que as \emph{features} com RSD dos QCs acima de um corte, como 20\%, são descartadas antes da análise estatística. Numa medição à parte, com cinco execuções do pipeline em R sobre o mesmo arquivo do HILIC positivo, em média 84 das 6.035 \emph{features} mudaram de lado do corte de 20\% entre duas execuções do SERRF, e 657 entre duas do LOESS. Nos outros três modos, as trocas foram de 6 a 15 \emph{features} com o SERRF e de 117 a 218 com o LOESS. A mesma análise, repetida sobre o mesmo arquivo, pode assim trabalhar com conjuntos diferentes de \emph{features}. Na ferramenta, duas execuções com a mesma entrada dão a mesma matriz.

A floresta da ferramenta passou por duas versões, e uma terceira foi testada. A primeira, com os parâmetros da biblioteca scikit-learn, corrigia bem menos que o R: no C18 positivo, deixava o RSD dos QCs fora da amostra em 21,55\%, contra 15,35\% do R. A causa era uma diferença de semântica. O tamanho mínimo para dividir um nó conta, no scikit-learn, as amostras distintas do sorteio com reposição e, no pacote ranger, usado pelo SERRF em R, os sorteios. Com sete QCs, cerca de 40\% das árvores da primeira versão não faziam divisão nenhuma. A floresta atual passa à árvore as amostras repetidas, como o ranger, e chega perto do R (Tabela~\ref{tab:ia-reais}): nos quatro modos ela corrige um pouco mais que ele, e em todos a diferença fica abaixo de um ponto percentual. Permitir ainda mais divisões por árvore afastaria mais, por exemplo 12,49\% contra 15,35\% do R no C18 positivo. A versão atual foi escolhida por ficar mais próxima do pipeline a que o laboratório está acostumado.

\par\medskip
\noindent\begin{minipage}{\linewidth}
  \tabelafixa{Mediana do RSD dos QCs fora da amostra (\%) nos quatro arquivos reais, com um lote}
  \label{tab:ia-reais}
  \centering
  \footnotesize
  \setlength{\tabcolsep}{3pt}
  \begin{tabular}{>{\raggedright\arraybackslash}p{2.0cm}>{\raggedright\arraybackslash}p{1.7cm}>{\raggedright\arraybackslash}p{1.9cm}>{\raggedright\arraybackslash}p{2.3cm}>{\raggedright\arraybackslash}p{1.9cm}>{\raggedright\arraybackslash}p{1.8cm}>{\raggedright\arraybackslash}p{2.3cm}}
    \hline
    \textbf{Modo} & \textbf{Sem correção} & \textbf{LOESS, ferramenta} & \textbf{LOESS, pipeline em R} & \textbf{SERRF, primeira versão} & \textbf{SERRF, versão atual} & \textbf{SERRF, pipeline em R} \\
    \hline
    HILIC positivo & 143,81 & 147,43 & 142,34 a 142,67 & 135,69 & 109,92 & 110,34 a 110,39 \\
    C18 positivo & 24,15 & 32,40 & 30,42 a 30,62 & 21,55 & 14,38 & 15,35 a 15,37 \\
    HILIC negativo & 67,18 & 59,20 & 59,57 a 59,80 & 66,08 & 61,45 & 61,63 a 62,10 \\
    C18 negativo & 15,50 & 19,32 & 18,76 a 19,19 & 15,45 & 15,12 & 15,17 a 15,37 \\
    \hline
  \end{tabular}
  \fonte{elaborado pelos autores.}
  \nota{um QC retirado por vez, remoção de valores discrepantes, semente 42 na ferramenta. O RSD do pipeline em R foi medido pelo mesmo procedimento, retirando um QC por vez, e não é o que o pipeline grava; ele foi executado três vezes, e a coluna dá o menor e o maior valor. Todas as medições no mesmo computador.}
\end{minipage}
\par\medskip

Pela medida fora da amostra, o LOESS da ferramenta fica pior que não corrigir em três modos e só fica à frente no HILIC negativo, e o do pipeline em R fica pior em dois, os C18.

O critério visual do laboratório, o agrupamento dos QCs no gráfico de componentes principais, foi medido pela razão entre a dispersão dos QCs e a das amostras nos dois primeiros componentes, que é tanto menor quanto mais agrupados estão os QCs. A versão atual do SERRF da ferramenta agrupa os QCs mais que o SERRF do pipeline nos quatro modos, com razões de 0,07 a 0,77, contra 0,08 a 0,79. Entre as correções do pipeline, porém, a que mais agrupa os QCs é o LOESS em três modos, e esse critério o elegeria no C18 positivo e no C18 negativo, onde, pela medida fora da amostra, ele deixa o RSD dos QCs 6,3 e 3,3 pontos percentuais acima de não corrigir. É a diferença entre avaliar a correção nos QCs que a ajustaram e avaliá-la em QCs retirados do ajuste.

No mesmo computador, o pipeline em R levou de 94 a 280 segundos por arquivo para aplicar o SERRF e o LOESS, e a ferramenta, pela linha de comando, de 78 a 242 segundos para aplicar os mesmos dois corretores, na mediana de três repetições. O ajuste do SERRF da ferramenta é mais lento que o do R (de 64 a 205 segundos, contra 36 a 143), e o do LOESS, mais rápido (de 11 a 33 segundos, contra 54 a 131). O pipeline em R usa metade dos núcleos do computador, como está configurado, e a ferramenta usa todos. Esses tempos cobrem só a etapa computacional; no fluxo anterior, a escolha do corretor e os cortes de \emph{features} eram feitos depois, à mão.

\subsection{O compromisso entre os QCs e as amostras}

A medida fora da amostra corrige o viés de avaliar a correção nos QCs que a ajustaram, mas continua olhando só para os QCs. Se uma correção reduz a variação dos QCs sem aproximar as amostras do valor verdadeiro, nenhuma medida feita nos QCs percebe. Nos arquivos reais esse valor não é conhecido. Nas matrizes sintéticas ele é, e a comparação pode ser feita por duas medidas. O erro relativo mediano é a mediana, sobre as células com valor verdadeiro positivo, da diferença absoluta entre o valor corrigido e o verdadeiro, dividida pelo verdadeiro. A fração de \emph{features} pioradas conta, para cada \emph{feature}, se o erro dela depois da correção ficou maior que sem correção, por qualquer quantidade. A Tabela~\ref{tab:ia-sinteticas} compara assim os corretores da ferramenta e os do pipeline em R, nas mesmas matrizes.

\par\medskip
\noindent\begin{minipage}{\linewidth}
  \tabelafixa{Erro relativo mediano (fração) contra a intensidade verdadeira e \emph{features} pioradas (\%) em matrizes sintéticas com deriva gradual}
  \label{tab:ia-sinteticas}
  \centering
  \footnotesize
  \setlength{\tabcolsep}{3pt}
  \begin{tabular}{>{\raggedright\arraybackslash}p{2.0cm}>{\raggedright\arraybackslash}p{1.7cm}>{\raggedright\arraybackslash}p{1.9cm}>{\raggedright\arraybackslash}p{2.3cm}>{\raggedright\arraybackslash}p{1.9cm}>{\raggedright\arraybackslash}p{1.8cm}>{\raggedright\arraybackslash}p{2.3cm}}
    \hline
    \textbf{Modo imitado} & \textbf{Sem correção} & \textbf{LOESS, ferramenta} & \textbf{LOESS, pipeline em R} & \textbf{SERRF, primeira versão} & \textbf{SERRF, versão atual} & \textbf{SERRF, pipeline em R} \\
    \hline
    HILIC positivo & 0,929 & 0,911; 17,6 & 0,898; 23,7 & 0,904; 37,1 & 0,889; 38,2 & 0,893; 38,1 \\
    C18 positivo & 0,228 & 0,190; 34,8 & 0,198; 36,2 & 0,169; 2,2 & 0,161; 7,8 & 0,162; 6,3 \\
    HILIC negativo & 0,749 & 0,782; 37,6 & 0,789; 39,4 & 0,602; 9,4 & 0,562; 10,0 & 0,569; 10,7 \\
    C18 negativo & 0,348 & 0,263; 24,3 & 0,265; 24,8 & 0,253; 2,7 & 0,197; 2,6 & 0,202; 2,9 \\
    \hline
  \end{tabular}
  \fonte{elaborado pelos autores.}
  \nota{médias de 20 matrizes por modo, razão entre deriva e ruído de 0,8, 300 \emph{features} e sete QCs; em cada célula, o erro e a porcentagem de \emph{features} pioradas em relação a não corrigir. O pipeline em R foi executado uma vez por matriz. O erro relativo pesa menos um valor que caiu do que um que subiu na mesma proporção. Em escala logarítmica, que trata igual um valor dobrado e um reduzido à metade, o LOESS da ferramenta aumenta o erro também nas matrizes que imitam o HILIC positivo (0,642 contra 0,499 sem correção) e piora 49,6\% das \emph{features}.}
\end{minipage}
\par\medskip

O SERRF reduz o erro em relação a não corrigir nos quatro modos, nas duas versões da ferramenta e no pipeline em R, e a versão atual fica muito perto do R. Ainda assim, a versão atual piora parte das \emph{features}, de 2,6\% a 38,2\%. Os dois LOESS pioram de 17,6\% a 39,4\% delas e, nas matrizes que imitam o HILIC negativo, aumentam o erro, justamente no modo em que o LOESS vence pela medida dos QCs nos arquivos reais. Em todas as células, o LOESS e o SERRF da ferramenta ficam perto dos do pipeline. Com o pipeline executado cinco vezes sobre cada uma dessas matrizes, a diferença típica, célula a célula, entre o valor corrigido pela ferramenta e pelo R foi de 1,8\% a 4,1\% no SERRF, de 1,5 a 1,9 vez a diferença entre duas execuções do próprio R (de 0,9\% a 2,4\%). No LOESS, que no R varia menos entre execuções (de 0,8\% a 1,5\%), foi de 1,8\% a 9,8\%, de 2,3 a 7,7 vezes. No C18 positivo, quanto mais divisões por árvore a floresta faz, mais \emph{features} pioram: 2,2\% na primeira versão, 7,8\% na atual e 12,7\% numa com mais divisões, com o erro praticamente igual. Nos outros três modos essa relação não apareceu. A piora também tem tamanho: em 20 matrizes por modo, entre as \emph{features} que pioraram com a versão atual do SERRF, o erro aumentou tipicamente de 12\% a 23\%, conforme o modo, e entre as que melhoraram, diminuiu de cerca de 20\% a 43\%.

A mesma comparação foi repetida, com 20 matrizes por modo, em outras quatro intensidades de deriva gradual e com um degrau no meio da corrida. Com deriva forte (razões de 1,5 e 3,0) e com o degrau, o SERRF e o LOESS reduziram o erro em relação a não corrigir nos quatro modos, e a fração de \emph{features} pioradas ficou abaixo de 18\%, salvo com o SERRF nas matrizes que imitam o HILIC positivo, de 14\% a 33\%. Com deriva fraca, o quadro se inverte. Na razão 0,2, o LOESS aumentou o erro nos quatro modos e o SERRF em três, e os dois pioraram de 36\% a 93\% das \emph{features}; na razão 0,4, o LOESS ainda aumentou o erro nos quatro, e o SERRF em dois. Nessas duas razões e no degrau, a floresta com mais divisões piorou mais \emph{features} que a primeira versão nos quatro modos. Com pouca deriva a remover, não corrigir foi, em geral, a melhor opção.

No degrau, também foi medido o efeito de declarar lotes. Em escala logarítmica, corrigir com os dois lotes separados exatamente no ponto do degrau reduziu o erro do SERRF nos quatro modos, e o do LOESS não mudou. Com os lotes que o detector propôs, presentes em 59 das 80 matrizes, o SERRF também errou menos nos quatro modos, e o LOESS errou mais em dois. Pelo erro relativo, nas matrizes que imitam o HILIC positivo, o SERRF com lotes parece piorar, pela mesma assimetria que a nota da Tabela~\ref{tab:ia-sinteticas} descreve.

A consequência para a escolha do corretor foi medida diretamente. Em 88 matrizes sintéticas com deriva gradual de intensidades variadas, o corretor de menor RSD dos QCs fora da amostra foi também o de menor erro contra a intensidade verdadeira em só 33 (37,5\%), ou em 38 (43,2\%) quando o erro é medido em escala logarítmica, que pesa igual um valor dobrado e um reduzido à metade. O SERRF foi o de menor erro em 35, não corrigir em 31 e o LOESS em 22. Escolher sempre o SERRF acertaria um pouco mais, 35, e a medida dos QCs nunca escolheu não corrigir. Ela também não o escolheu onde a correção piora as amostras. Na razão 0,2, não corrigir teve o menor erro contra a verdade nas 20 matrizes de três dos modos, e a medida dos QCs o escolheu em só 2 dessas 60, sem se abster em nenhuma; nas outras 58, escolheu uma correção, quase sempre o SERRF (55). Na razão 0,4, não o escolheu em nenhuma das 80 matrizes, e se absteve em duas. Escolher pela medida dos QCs, portanto, não garante a melhor correção para as amostras, e o fluxo anterior do laboratório, que também escolhe olhando só os QCs, está sujeito ao mesmo limite. A literatura aponta o mesmo limite: avaliar a correção só pelos QCs que a ajustaram tende a dar um resultado otimista e não mede o efeito sobre as amostras \cite{martens2023instrumental}, e há casos em que o RSD dos QCs cai sem que amostras de referência independentes melhorem \cite{brunius2016large}.

Diante disso, a ferramenta inclui não corrigir entre as opções, a medição completa pode se abster e a caixa ``Automático'' pode ser desmarcada: a decisão final fica com quem conhece o estudo. Uma política treinada pela intensidade verdadeira das matrizes sintéticas foi construída e avaliada no gerador, mas ficou desligada, porque nada garante que o que vale para o gerador valha para dados reais.

\subsection{Aprendizado de máquina nas análises}

Nas análises estatísticas, os modelos supervisionados vêm sempre acompanhados de validação. O PLS-DA só devolve escores junto com o \(Q^2\), a fração da variação dos rótulos que o modelo prevê, estimada por validação cruzada estratificada em sete partes, e com o \emph{p}-valor de um teste de permutação com 200 sorteios, como recomendam Westerhuis e colaboradores \cite{westerhuis2008assessment}. Os dois indicadores não têm a mesma calibração. Em 100 repetições com rótulos de grupo sorteados ao acaso, o \emph{p}-valor de permutação ficou acima de 0,05 em 94 a 95 delas em cada modo, como esperado, mas o \(Q^2\) saiu positivo em até 14, chegando a 0,47. Um \(Q^2\) positivo, sozinho, não indica que a separação seja real. Os classificadores por floresta aleatória e por máquina de vetores de suporte são avaliados da mesma forma, com acurácia balanceada, nível do acaso e \emph{p}-valor de permutação. A exigência importa porque, com 61 amostras e até 6.035 \emph{features} por corrida, um classificador pode separar perfeitamente os grupos dentro da amostra mesmo sem diferença real entre eles.

A escolha do teste univariado também é automática. Para cada par de grupos, um único teste de Shapiro-Wilk \cite{shapiro1965analysis} decide, para todas as \emph{features} de uma vez, entre o teste \emph{t} de Welch e o de Mann-Whitney, para que a correção de Benjamini-Hochberg seja aplicada a \emph{p}-valores de um mesmo teste, e o teste escolhido fica declarado no gráfico de vulcão, com o \emph{p}-valor do teste de normalidade. O teste é aplicado aos desvios de cada valor em relação à média do seu grupo, limitados a 5.000 valores tomados em posições igualmente espaçadas, para que todas as amostras dos dois grupos entrem na decisão. Com milhares de valores, o Shapiro-Wilk rejeita a normalidade diante de desvios pequenos, e \emph{features} de escalas diferentes formam uma mistura que não é normal mesmo quando cada uma é. Em simulação com dados normais dentro de cada grupo, ele rejeitou a normalidade em 3 de 40 conjuntos com escalas iguais, em 5 de 40 com transformação logarítmica e autoescalonamento e em 40 de 40 com escalas diferentes e sem escalonamento. A escolha do teste depende, portanto, do pré-processamento. Nestas corridas, com transformação logarítmica e autoescalonamento, o teste de Mann-Whitney foi o escolhido em todos os contrastes, com os três corretores da ferramenta e com as matrizes corrigidas pelo pipeline em R, e nenhuma \emph{feature} ficou abaixo de \emph{q} = 0,05 em nenhum dos 24 contrastes de cada um (seis pares de grupos nos quatro modos). Nesse resultado, a escolha do corretor não muda a conclusão.
